\section{Finite certificates for unbounded behaviour}
\label{sec:closure}

\subsection{The shape of the problem}

A large class of obligations has the form: \emph{for every input word, every
number of iterations, every reachable state, some observation is zero}. Program
equivalence, fold correctness, loop invariants and accumulator specifications
all land here. The universal quantifier ranges over an infinite set, so no
amount of testing settles it --- but the state space that matters is often
finite-dimensional, and that is what makes a finite certificate possible.

\S\ref{sec:structural} approached this by \emph{guessing}: propose an invariant
template, solve for its coefficients, check the base and step identities. That
works, and Proposition~\ref{prop:rec-complete} shows it is complete on its
fragment. But it has two defects that the same section had to admit.

The first is that the implemented certificate language requires a
\emph{conserved} quantity, $I(T(s))=I(s)$. That is strictly too narrow. Take
\[
 x'=2x+a,\qquad y'=2y+2a,\qquad (x_0,y_0)=(0,0),\qquad a\in\{0,1\},
\]
and the observation $q=2x-y$. After every input word $q$ is zero --- yet it is
not conserved: the rule is $q'=2q$, not $q'=q$. A worker that can only express
conservation cannot state this invariant at all, let alone prove it.

The second is the bilinearity trap. For a \emph{family} of invariants the step
condition is
\[
 I_i(T(\vec x))=\sum_j h_{ij}(\vec x)\,I_j(\vec x),
\]
and \S\ref{sec:structural} warns that discovering the $I_j$ and the multipliers
$h_{ij}$ together is a bilinear search which must not be described as linear
algebra, recommending only that one ``fix one side at a time''.

This section is what fixing one side at a time looks like when it is done
constructively. The idea is to stop guessing and start from the \emph{target}:
take the observation one actually cares about, close it under pullback along
the transition until the closure stabilises, and then solve for multipliers
over the resulting \emph{fixed} family --- which is an ordinary linear problem.
The closure either terminates with a certificate valid for every input, or it
produces a concrete input on which the observation is nonzero. There is no
template to choose and no ansatz degree to tune.

\subsection{Observable-space closure}
\label{sec:obs-closure}
\status{Implemented and tested in Python by several independent efforts.}

\subsubsection*{The fragment}

Let $A_a\in\Q^{D\times D}$ for each letter $a\in\Sigma$, let $s_0\in\Q^{D}$, and
let $q\in\Q^{1\times D}$. The goal is
\begin{equation}\label{eq:linear-goal}
 \forall w\in\Sigma^{*},\qquad q\,A_w\,s_0=0,
 \qquad A_{a_1\cdots a_t}=A_{a_t}\cdots A_{a_1}.
\end{equation}
\emph{All} words are permitted. A restricted input language needs either an
explicit product with a supplied finite automaton or a separate proof that the
restriction is respected; it must never be silently assumed.

Affine updates are included by homogenisation,
\[
 s'=M_as+b_a
 \quad\leadsto\quad
 \begin{bmatrix}s'\\1\end{bmatrix}
 =\begin{bmatrix}M_a&b_a\\0&1\end{bmatrix}\begin{bmatrix}s\\1\end{bmatrix},
\]
and an affine target likewise. Rational constants stay rational: no
floating-point rank test is needed, and none would be appropriate.

Equivalence of two implementations is the same problem. Form the block system
\[
 A_a=\begin{bmatrix}A^{(1)}_a&0\\0&A^{(2)}_a\end{bmatrix},\qquad
 q=\begin{bmatrix}o_1&-o_2\end{bmatrix},
\]
and a single zero-observation goal compares every output. Equal output requires
neither equal internal states, nor equal state dimensions, nor a shared
coordinate system --- which is exactly what makes this useful for proving two
differently written folds equal.

\subsubsection*{The algorithm}

Start from the target row and close it under right multiplication.

\begin{lstlisting}
queue := [(q, empty_word)]
basis := []
while queue is not empty:
    (r, w) := pop_front(queue)
    if r * initial != 0:
        replay w in the original matrices
        return Counterexample(w)
    if r lies in span(basis):
        continue
    append r to basis
    for each action a:
        push_back(queue, (r * A[a], prepend(a, w)))
return the closure coefficients for the completed basis
\end{lstlisting}

Dependent rows need not be expanded: multiplying their known linear combination
by $A_a$ produces the corresponding combination of rows already scheduled. Each
queued row carries a concrete word, and because a row representing $qA_w$
becomes $qA_wA_a$ under right multiplication, that word grows on the \emph{left}
--- the successor of $w$ is $aw$, not $wa$. Getting this backwards yields a
``counterexample'' that does not reproduce.

\subsubsection*{A certificate that needs no rank computation}

The search returns a matrix $G\in\Q^{r\times D}$, a row $t\in\Q^{1\times r}$,
and matrices $C_a\in\Q^{r\times r}$ with
\begin{align}
 q&=tG, \label{eq:lin-target}\\
 GA_a&=C_aG \qquad (a\in\Sigma), \label{eq:lin-closure}\\
 Gs_0&=0. \label{eq:lin-init}
\end{align}

\begin{theorem}[Observable certificate soundness]
\label{thm:linear-sound}
If \eqref{eq:lin-target}--\eqref{eq:lin-init} hold then \eqref{eq:linear-goal}
holds.
\end{theorem}
\begin{proof}
The invariant is $Gs=0$. If it holds before action $a$ then
$G(A_as)=(GA_a)s=(C_aG)s=C_a(Gs)=0$. It holds initially by
\eqref{eq:lin-init}, hence after every finite word by induction. Finally
$qs=tGs=0$ by \eqref{eq:lin-target}.
\end{proof}

The checker verifies three finite matrix identities and nothing else. It does
\emph{not} check that the rows of $G$ are independent, that they came from any
particular search, or that $r$ is minimal. Those properties explain the
algorithm's efficiency; they are irrelevant to soundness, and a certificate
carrying redundant rows is still valid. This is the cleanest search/check
separation anywhere in this document: the producer may be clever, the checker
does linear algebra.

Conversely \eqref{eq:lin-target} is indispensable. An invariant basis without
the target-membership equation may prove something perfectly true and entirely
irrelevant.

\subsubsection*{Termination, completeness, and a sharp witness bound}

\begin{theorem}[Finite-dimensional coverage]
\label{thm:linear-complete}
Without operational cutoffs, observable closure terminates and \emph{decides}
\eqref{eq:linear-goal}. If the goal is false and $q\ne0$, a distinguishing word
of length at most $D-1$ exists.
\end{theorem}
\begin{proof}
At most $D$ independent rows can be admitted, and each admitted row creates
$|\Sigma|$ successors, so the queue holds at most $1+|\Sigma|D$ entries over the
whole run. On normal termination every basis row is closed under right
multiplication by every action, so induction places every $qA_w$ in the final
row space; if all basis rows annihilate $s_0$ then so does every observation,
and otherwise a tested row supplies a concrete word.

For the bound, let $V_i$ be the span of the target rows arising from words of
length at most $i$. If $V_i=V_{i+1}$ the space is closed under each action and
the chain has stabilised. Starting from the one-dimensional span of a nonzero
$q$, at most $D-1$ strict increases are possible, so $V_{D-1}$ contains every
target row; if its evaluation at $s_0$ is not identically zero, one of its
generating words of length at most $D-1$ has nonzero evaluation.
\end{proof}

The bound is attained, and the family that attains it is the best possible
argument against testing. Take a shift matrix with ones just above the
diagonal, initial state $e_D$, and target $e_1^{\mathsf T}$: every output before
step $D-1$ is zero, and the next one is one. \emph{Testing only the first
several inputs would classify these systems as correct.} What tells us when a
conclusion is justified is the finite-dimensional argument, not the number of
samples --- which is the same lesson as \S\ref{sec:structural}'s interpolation
trap, in a setting where the trap is invisible to any finite test.

With a maintained echelon basis, row multiplication and membership reduction
cost $O(D^2)$ field operations per candidate, so an $O(|\Sigma|D^3)$
implementation exists. Bit complexity additionally depends on exact coefficient
growth, and the prototypes recompute ranks symbolically rather than maintaining
the incremental representation, so the cubic figure is a production target and
not a measured property of the delivered code.

\subsubsection*{Beyond a finite alphabet, carefully}

Suppose an update depends polynomially on a fresh scalar input,
$A(u)=\sum_{j=0}^{d}u^jM_j$. A sufficient universal certificate is
\[
 GM_j=C_jG \qquad (0\le j\le d),
\]
since then $GA(u)=\bigl(\sum_j u^jC_j\bigr)G$ for every $u$. Over an infinite
field, an output polynomial vanishing for all inputs forces all its
coefficients to vanish, so coefficient closure is complete for this precise
polynomial-input, linear-state fragment.

\begin{remark}[A coefficient matrix is not an input symbol]
If a coefficient row fails to annihilate the initial state, its index sequence
is evidence of a nonzero output \emph{polynomial} --- not a counterexample.
To produce an actual counterexample, instantiate that polynomial on an exact
grid whose coordinate sizes exceed the respective degrees, then replay the
chosen rational inputs. Presenting the coefficient indices to a user as though
they were inputs would be a semantic bug of exactly the kind
Invariant~\ref{inv:tuple} exists to prevent.
\end{remark}

\subsection{Target-generated ideal closure}
\label{sec:ideal-closure}
\status{Implemented and tested in Python by several independent efforts.}

Linear closure decides linear observations of linear systems. The polynomial
case replaces the row space by an \emph{ideal}: pull the target and its
auxiliary generators back through the updates and close the ideal they generate,
rather than closing a span.

The certificate is correspondingly richer. Where the linear case emits closure
matrices $C_a$, the polynomial case emits \emph{polynomial multipliers}: each
pulled-back generator is expressed as a combination of the current generators
with polynomial coefficients. The checker verifies those identities directly by
exact polynomial arithmetic. In particular it does not have to trust a
Gr\"obner-basis computation, or even to know that one occurred --- the
multiplier identities are the certificate, and a basis computation is merely
how the producer found them.

This is the constructive answer to the bilinearity trap. The invariants are not
guessed jointly with their multipliers: the family is \emph{constructed} by
closure, and only then are multipliers solved for over that fixed family, which
is linear.

\subsection{What a counterexample here actually refutes}

The closure workers return concrete negative results, and those are worth more
than the usual \unknown{}: a distinguishing word for the linear case, a
concrete orbit point or execution for the polynomial case. Each is validated
independently --- the word is replayed in the \emph{original} matrices, not in
the closure representation.

But the scope statement matters as much as the object.

\begin{remark}[Scope of a distinguishing word]
A replayed word refutes \eqref{eq:linear-goal} \emph{as modelled}. It refutes
the source obligation only to the extent that the modelling is faithful: that
the matrices really are the source updates, that the alphabet really is the
input type, and --- the one most easily forgotten --- that the source admits
\emph{all} words. If the source restricts its inputs and the model did not,
a distinguishing word may be unreachable in the source, and the correct report
is that the model was wrong, not that the program is.
\end{remark}

This places closure counterexamples in the \emph{countermodel replayed} row of
the status lattice (\S\ref{sec:lattice}): they refute the modelled statement for
this system, subject to the recorded assumptions --- which is a strong and
useful result, and is not the same as refuting the Lean goal.
